\metatitle{Total positivity}
\metadescription{Total positivity and total nonnegativity as organizing
principles for matrices, planar networks, Pólya frequency sequences,
positroids, and real-rootedness.}
\metakeywords{total positivity,total nonnegativity,TNN matrix,planar network,
Pólya frequency sequence,positroid}

\section[totalPositivity]{Total positivity and total nonnegativity}

Total positivity is the study of matrices and kernels whose minors have a
fixed sign.  A matrix is \hyperref[totallyNonnegative]{totally nonnegative}
(TNN) if all minors are nonnegative, and
\hyperref[totallyPositiveMatrix]{totally positive} if all minors are
positive.  See \cite{Karlin1968,FallatJohnson2011} for general background.

This page is a navigation point for the versions that occur most often in
algebraic combinatorics: planar networks, Pólya frequency sequences, real
zeros, and totally nonnegative Grassmannians.

\begin{example}
The matrix
\[
  \begin{pmatrix}
  1 & 1\\
  1 & 2
  \end{pmatrix}
\]
is totally positive: all entries are positive, and its determinant is \(1\).
For larger matrices, checking all minors is usually inefficient, so one uses
structure such as planar networks, Toeplitz form, or factorization.
\end{example}

\subsection[totalPositivityNetworks]{Planar networks}

The \hyperref[lgvLemma]{Lindström--Gessel--Viennot lemma} explains why planar
networks produce TNN matrices.  If a matrix entry counts weighted paths from
a source to a sink in a compatible planar network, then each minor counts
families of nonintersecting paths and is therefore nonnegative.

\svgimg[width=0.82\textwidth]{svg-images/total-positivity-network.svg}{
A planar network in which nonintersecting path families give nonnegative
minors.}

This mechanism appears in:
\begin{itemize}
\item \hyperref[tnnPathMatrix]{path matrices} for total positivity;
\item Jacobi--Trudi and flagged Jacobi--Trudi determinants;
\item \hyperref[positroidNetworks]{planar networks and plabic graphs} for
\hyperref[positroid]{positroids};
\item network proofs of real-rootedness and interlacing.
\end{itemize}

\subsection[totalPositivityPF]{Pólya frequency sequences}

A sequence \(a=(a_0,a_1,\dotsc)\) is a
\hyperref[polyaFrequencySequenceDefinition]{Pólya frequency sequence} if its
Toeplitz matrix \(T(a)_{i,j}=a_{i-j}\) is TNN.  The
\hyperref[aissenSchoenbergWhitney]{Aissen--Schoenberg--Whitney theorem}
identifies finite Pólya frequency sequences with polynomials whose zeros are
real and nonpositive.

Thus total positivity gives one of the main bridges between coefficient
sequences and \hyperref[realRootedPolynomial]{real-rooted polynomials}.
It also explains closure properties: products of TNN matrices are TNN, so
convolutions of Pólya frequency sequences remain Pólya frequency sequences.

\subsection[totalPositivityGrassmannians]{Grassmannians and positroids}

For a \(k\times n\) matrix \(A\), the maximal minors \(\Delta_I(A)\) are the
\hyperref[pluckerCoordinate]{Plücker coordinates} of a point of the
\hyperref[grassmannianVarieties]{Grassmannian} \(\mathrm{Gr}(k,n)\).
The totally nonnegative Grassmannian consists of points with all Plücker
coordinates nonnegative.  Its cells are indexed by
\hyperref[positroid]{positroids}, introduced by \name{Alexander Postnikov}
\cite{Postnikov2006x}.

This is a different but closely related face of total positivity: instead of
asking all minors of a square matrix to be nonnegative, one fixes the maximal
minors that define a Grassmannian point.

\subsection[totalPositivityRealRootedness]{Real-rootedness}

Total positivity frequently packages root-location statements.  The
\hyperref[polyaFrequency]{Pólya frequency page} treats Toeplitz total
positivity and real-rootedness of coefficient sequences.  The
\hyperref[interlacingPolynomials]{interlacing page} uses total
nonnegativity to encode pairs and families of interlacing polynomials.  The
\hyperref[stablePolynomial]{stable-polynomial page} gives a multivariate
framework that often implies the same univariate real-rootedness results
after specialization.

In practice, the useful question is often not "are all minors nonnegative?"
but "which structured matrix has the desired minors?"  Once the matrix is
identified as a path matrix, Toeplitz matrix, or product of elementary TNN
matrices, the desired positivity follows from closure properties.
